\documentclass[a4paper,11pt]{article}
\usepackage[utf8]{inputenc}
\usepackage[T1]{fontenc}
\usepackage[french]{babel}
\usepackage{graphicx}
\usepackage{amsmath}
\usepackage{amsthm}
\theoremstyle{definition}
\usepackage{geometry}
\usepackage{fancyhdr}
\usepackage{hyperref}
\usepackage[usenames,dvipsnames,svgnames,table]{xcolor}
\usepackage{enumerate}
\usepackage{tcolorbox}
\tcbuselibrary{skins,breakable}
\usepackage{tikz}
\usepackage{pgfplots}
\pgfplotsset{compat=1.18}

\geometry{a4paper, scale=0.8}

%% ─── Header / footer ─────────────────────────────────────────────
\pagestyle{fancy}
\fancyhf{}
\fancyhead[L]{2025-2026}
\fancyhead[C]{\textbf{\courseName}}
\fancyhead[R]{\large GEII -- BUT2}
\fancyfoot[C]{\thepage}
\fancyfoot[L]{\includegraphics[height=3em]{../../common/logos_couleurs/IUTA.png}}
% \fancyfoot[R]{\href{mailto:prenom.nom@univ-smb.fr}{prenom.nom@univ-smb.fr}}

%% ─── Course metadata ─────────────────────────────────────────────
\newcommand{\courseName}{Outils Mathématiques et Logiciel --- OML3}
\newcommand{\chapterTitle}{TD 4 --- Transformée de Fourier}

%% ─── Custom title ────────────────────────────────────────────────
\makeatletter
\renewcommand\maketitle{{\raggedright\begin{center}
    {\Huge\bfseries\@title}
\end{center}}}
\makeatother

%% ─── Solution toggle ─────────────────────────────────────────────
%% Compile with \showsolutionstrue for the teacher version.
\newif\ifshowsolutions
\showsolutionsfalse

\ifshowsolutions
    \newtcolorbox{solution}{
        colback=white, colframe=white,coltext=red,
        breakable, title=\textbf{Solution :}
    }
\else
    \NewEnviron{solution}{}{}
\fi

%% ─── Theorem / exercise environments ─────────────────────────────
\newtheorem{definition}{Définition}[section]
\newtheorem{proposition}{Propriété}[section]
\newtheorem{theorem}{Théorème}[section]
\newtheorem{exemple}{Exemple}[section]
\newtheorem{remark}{Remarque}[section]

%% Exercice : affiche "Exercice N (titre)." et referme aussitôt le
%% paragraphe avec \par. Volontairement PAS un \newtheorem : amsthm
%% reporte l'affichage du titre via \everypar, qui ne se déclenche
%% qu'au démarrage d'un paragraphe. Si le corps commence directement
%% par une liste (\begin{enumerate}), aucun paragraphe ne démarre
%% avant le premier \item, et le titre se retrouve collé au "1." de
%% la liste -- même avec une ligne vide dans la source. Cet
%% environnement affiche le titre immédiatement, donc ne dépend pas
%% de ce mécanisme et fonctionne quel que soit le contenu qui suit.
\newcounter{exercice}   % global counter — do not add [section]
\newenvironment{exercice}[1][]{%
  \refstepcounter{exercice}%
  \par\noindent\textbf{Exercice \theexercice}\if\relax\detokenize{#1}\relax\else\ (#1)\fi.\par\noindent\ignorespaces
}{\par}

%% ─── Section numbering ───────────────────────────────────────────
\renewcommand{\thesection}{{\arabic{section}}}

%% ─── Helpers ─────────────────────────────────────────────────────
\definecolor{niceblue}{RGB}{0,128,172}
\newcommand{\highlight}[1]{{\color{niceblue}\emph{#1}}}

%% ─── Peigne de Dirac (Ш) ─────────────────────────────────────────
%% Pas de glyphe cyrillique dans les fontes Computer Modern : on le
%% dessine (trois barres + base), indépendant de toute fonte externe.
\newcommand{\Sha}{%
  \mathbin{\vcenter{\hbox{%
    \begin{tikzpicture}[line width=0.55pt, y=0.62em, x=0.62em]
      \draw (0,0) -- (0,1);
      \draw (0.5,0) -- (0.5,1);
      \draw (1,0) -- (1,1);
      \draw (0,0) -- (1,0);
    \end{tikzpicture}%
  }}%
  }%
}

\usepackage{amsmath,amsfonts,bm}

%% ─── Operators ───────────────────────────────────────────────────
\DeclareMathOperator*{\argmax}{arg\,max}
\DeclareMathOperator*{\argmin}{arg\,min}
\DeclareMathOperator{\MSE}{MSE}
\DeclareMathOperator{\sign}{sign}
\DeclareMathOperator{\sinc}{sinc}
\DeclareMathOperator{\tr}{tr}

%% ─── Number sets ─────────────────────────────────────────────────
\def\R{{\mathbb{R}}}
\def\C{{\mathbb{C}}}
\def\N{{\mathbb{N}}}
\def\Z{{\mathbb{Z}}}

%% ─── Infinitesimal and units ────────────────────────────────────
\def\d{{\rm d}}
\def\Hz{{\mbox{Hz}}}
\def\m{{\mbox{m}}}

%% ─── Complex / signal ───────────────────────────────────────────
\DeclareMathOperator{\complexj}{\mathsf{j}}   % imaginary unit

%% ─── Probability ─────────────────────────────────────────────────
\newcommand{\Esp}[2]{\mathbb{E}_{#1}\!\left[#2\right]}
\renewcommand{\Pr}{\mathbb{P}}

%% ─── Linear algebra ──────────────────────────────────────────────
\newcommand{\abs}[1]{\left|{#1}\right|}
\newcommand{\norm}[1]{\left\|#1\right\|}
\newcommand{\inner}[2]{\left\langle#1,\,#2\right\rangle}

%% ─── Logic ───────────────────────────────────────────────────────
\newcommand{\then}{\Rightarrow}
\newcommand{\iif}{\Leftrightarrow}

%% ─── Draft annotation ────────────────────────────────────────────
\newcommand{\red}[1]{{\color{red}#1}}

%% ─── Shared figure color (used by common/figures/*/*.tex bodies) ──
\definecolor{niceblue}{RGB}{0,128,172}

%% ─── Rappel header (styled subtitle for reference-card asides) ───
%% Requires \definecolor{niceblue}{...} to be defined by the including document.
\newcommand{\rappelhead}[1]{%
    \bigskip
    {\color{niceblue}\bfseries #1}\\[-0.4em]
    {\color{niceblue!35}\rule{\linewidth}{0.6pt}}
    \smallskip

}


\title{\chapterTitle}
\author{}
\date{}

%% ─────────────────────────────────────────────────────────────────
\begin{document}
\fontfamily{bch}\selectfont
\fontsize{12pt}{14pt}\selectfont

\maketitle
\thispagestyle{fancy}
\setcounter{exercice}{0}


%% ── Ex 1 : PROVE ─────────────────────────────────────────────────
%% Purpose: make students re-derive a theorem stated in the CM.
\begin{exercice}
    \label{ex:dualite-conv}

    \textit{Pour deux signaux $x,y:\R\to\R$ (des fonctions du temps $t$) tels que $x(t)\overset{\mathcal F}{\longleftrightarrow}X(f)$ et $y(t)\overset{\mathcal F}{\longleftrightarrow}Y(f)$ :
    $$
        X(f) = \int_{-\infty}^{+\infty} x(t)\,e^{-2\complexj\pi ft}\,\d t, \qquad\qquad (x*y)(t) = \int_{-\infty}^{+\infty} x(\tau)\,y(t-\tau)\,\d\tau.
    $$
    On admettra que, dans les calculs de cet exercice, l'ordre d'intégration d'une intégrale double peut être échangé, et qu'un changement de variable affine (par exemple $u=t-\tau$, à $\tau$ fixé) ne modifie pas des bornes $-\infty$ à $+\infty$.
    }

    \bigskip
    Le cours affirme, sans démonstration, que
    $$
    (x*y)(t) \overset{\mathcal F}{\longleftrightarrow} X(f)\,Y(f)
    $$
    (propriété de dualité produit/convolution). En utilisant les rappels ci-dessus, le démontrer.

    \medskip
    \noindent\textbf{Indication.} Pour tous $u,t\in\R$, $t = t+u-u$.
    \begin{solution}
        Par définition de la TF et du produit de convolution,
        $$
            \mathcal F[x*y](f) = \int_{-\infty}^{+\infty}\left(\int_{-\infty}^{+\infty} x(\tau)\,y(t-\tau)\,\d\tau\right) e^{-2\complexj\pi ft}\,\d t.
        $$
        En échangeant l'ordre d'intégration :
        $$
            \mathcal F[x*y](f) = \int_{-\infty}^{+\infty} x(\tau) \left(\int_{-\infty}^{+\infty} y(t-\tau)\,e^{-2\complexj\pi ft}\,\d t\right) \d\tau.
        $$
        Dans l'intégrale intérieure, on pose $u=t-\tau$ (donc $t=u+\tau$, $\d u=\d t$) : quand $t$ décrit $\R$ tout entier, $u$ aussi (pour tout $u\in\R$, $t=u+\tau$ convient), donc les bornes ne changent pas :
        $$
            \int_{-\infty}^{+\infty} y(t-\tau)\,e^{-2\complexj\pi ft}\,\d t = \int_{-\infty}^{+\infty} y(u)\,e^{-2\complexj\pi f(u+\tau)}\,\d u = e^{-2\complexj\pi f\tau}\int_{-\infty}^{+\infty} y(u)\,e^{-2\complexj\pi fu}\,\d u = e^{-2\complexj\pi f\tau}\,Y(f).
        $$
        D'où
        $$
            \mathcal F[x*y](f) = \int_{-\infty}^{+\infty} x(\tau)\,e^{-2\complexj\pi f\tau}\,Y(f)\,\d\tau = Y(f)\int_{-\infty}^{+\infty} x(\tau)\,e^{-2\complexj\pi f\tau}\,\d\tau = Y(f)\,X(f).
        $$
    \end{solution}
\end{exercice}


% %% ── Ex 2 : APPLY ─────────────────────────────────────────────────
% \begin{exercice}[Filtrer un signal échantillonné]
%     Soient $f_1,\dots,f_K\in\R$ des fréquences fixées et
%     $$
%         h(t) = \sum_{k=1}^K e^{-2\complexj\pi f_kt}.
%     $$
%     Soient $t_1,\dots,t_N\in\R$ des instants fixés, $a_1,\dots,a_N\in\R$ des coefficients, et
%     $$
%         x(t) = \sum_{n=1}^N a_n\,\delta(t-t_n)
%     $$
%     (un signal \highlight{échantillonné} : une somme pondérée d'impulsions).
%     \begin{enumerate}[\bfseries 1.]
%         \item Calculer $H=\mathcal F[h]$. On rappelle que
%         $\mathcal F\big[e^{2\complexj\pi f_0t}\big](f) = \delta(f-f_0)$ pour tout $f_0\in\R$.
%         \begin{solution}
%             Avec $f_0=-f_k$ : $\mathcal F\big[e^{-2\complexj\pi f_kt}\big](f) = \delta(f+f_k)$. Par linéarité,
%             $$
%                 H(f) = \sum_{k=1}^K \delta(f+f_k).
%             $$
%         \end{solution}

%         \item L'objet $x*h$ est-il un nombre ou une fonction de $t$ ? Que représente $u$ dans la définition $$(x*h)(t) = \displaystyle\int_{-\infty}^{+\infty} x(u)\,h(t-u)\,\d u.$$ Est-ce que changer le nom de la variable $u$ en $\tau$ change $x * h$ ?
%         \begin{solution}
%             Une fonction de $t$ : la variable $\tau$ disparaît une fois l'intégration effectuée, mais $t$ reste présent (dans $h(t-\tau)$) et n'est pas intégré. Chaque valeur de $t$ donne un nombre, différent en général d'une valeur de $t$ à l'autre : l'ensemble de ces nombres est une fonction de $t$, notée $x*h$.
%         \end{solution}

%         \item Un filtre $T$ (qui à un signal d'entrée $x$ associe un signal de sortie $T(x)$) est dit \highlight{linéaire} si $T(\alpha x_1+\beta x_2) = \alpha\,T(x_1)+\beta\,T(x_2)$ pour tous signaux $x_1,x_2$ et tous $\alpha,\beta\in\R$. Montrer que $T:x\mapsto x*h$ est linéaire, c'est-à-dire que $(\alpha x_1+\beta x_2)*h = \alpha(x_1*h)+\beta(x_2*h)$.
%         \begin{solution}
%             $$
%                 \big[(\alpha x_1+\beta x_2)*h\big](t) = \int_{-\infty}^{+\infty} \big(\alpha x_1(\tau)+\beta x_2(\tau)\big)\,h(t-\tau)\,\d\tau.
%             $$
%             Par linéarité de l'intégrale, cette expression se sépare en
%             $$
%                 \alpha\int_{-\infty}^{+\infty} x_1(\tau)\,h(t-\tau)\,\d\tau + \beta\int_{-\infty}^{+\infty} x_2(\tau)\,h(t-\tau)\,\d\tau = \alpha\,(x_1*h)(t) + \beta\,(x_2*h)(t).
%             $$
%         \end{solution}

%         \item $T$ est dite de plus \highlight{invariante dans le temps} si, pour tout $t_0\in\R$, retarder l'entrée de $t_0$ (remplacer $x$ par $t\mapsto x(t-t_0)$) revient à retarder la sortie de la même durée (obtenir $t\mapsto y(t-t_0)$, où $y=T(x)$). Une transformation à la fois linéaire et invariante dans le temps est dite \highlight{LTI}.

%         On pose $x_0(t)=x(t-t_0)$ et $y=x*h$. On va montrer, en quatre petits pas, que $(x_0*h)(t) = y(t-t_0)$ pour tout $t$ -- c'est-à-dire que $T:x\mapsto x*h$ est bien invariante dans le temps.
%         \begin{enumerate}[\bfseries a.]
%             \item Écrire $(x_0*h)(t)$ comme une intégrale, à partir de la définition de la convolution.
%             \begin{solution}
%                 $$
%                     (x_0*h)(t) = \int_{-\infty}^{+\infty} x_0(\tau)\,h(t-\tau)\,\d\tau.
%                 $$
%             \end{solution}
%             \item Remplacer $x_0(\tau)$ par sa définition dans cette intégrale.
%             \begin{solution}
%                 $$
%                     (x_0*h)(t) = \int_{-\infty}^{+\infty} x(\tau-t_0)\,h(t-\tau)\,\d\tau.
%                 $$
%             \end{solution}
%             \item Poser $u=\tau-t_0$ (donc $\tau=u+t_0$, $\d u=\d\tau$) et réécrire l'intégrale en fonction de $u$.
%             \begin{solution}
%                 $$
%                     (x_0*h)(t) = \int_{-\infty}^{+\infty} x(u)\,h\big((t-t_0)-u\big)\,\d u.
%                 $$
%             \end{solution}
%             \item Reconnaître, dans le résultat, l'expression de $y$ en un certain point. Conclure.
%             \begin{solution}
%                 L'intégrale obtenue est exactement $(x*h)$ évaluée en $t-t_0$, c'est-à-dire $y(t-t_0)$. Donc $(x_0*h)(t)=y(t-t_0)$ pour tout $t$ : $T$ est invariante dans le temps. Avec la question précédente, $T:x\mapsto x*h$ est LTI, de \highlight{réponse impulsionnelle} $h$.
%             \end{solution}
%         \end{enumerate}

%         \item Calculer $X=\mathcal F[x]$. On rappelle que $\mathcal F[\delta(\cdot-t_0)](f) = e^{-2\complexj\pi ft_0}$.
%         \begin{solution}
%             Par linéarité,
%             $$
%                 X(f) = \sum_{n=1}^N a_n\,e^{-2\complexj\pi ft_n}.
%             $$
%         \end{solution}

%         \item Le théorème produit-convolution énonce que $\mathcal F[u*v] = \mathcal F[u]\,\mathcal F[v]$ pour tous signaux $u,v$. En déduire $\mathcal F[y]$, où $y=h*x$, puis $y(t)$. On rappelle que $\delta(f+f_k)\,X(f) = X(-f_k)\,\delta(f+f_k)$ (multiplier une fonction par une impulsion décalée revient à multiplier l'impulsion par la valeur de la fonction en ce point).
%         \begin{solution}
%             $$
%                 \mathcal F[y](f) = H(f)X(f) = \left(\sum_{k=1}^K \delta(f+f_k)\right) X(f) = \sum_{k=1}^K X(-f_k)\,\delta(f+f_k), \qquad X(-f_k) = \sum_{n=1}^N a_n\,e^{2\complexj\pi f_kt_n}.
%             $$
%             En repassant en temporel (même calcul qu'à la question 1, dans l'autre sens) :
%             $$
%                 y(t) = \sum_{k=1}^K X(-f_k)\,e^{-2\complexj\pi f_kt}.
%             $$
%             Convoler $x$ par $h$ ne conserve donc, du signal échantillonné $x$, que les $K$ valeurs de son spectre aux fréquences $-f_1,\dots,-f_K$ : $h$ agit comme un filtre qui ne laisse passer qu'un nombre fini de fréquences.
%         \end{solution}

%         \item \textbf{Conclusion.} Que devient le résultat de la question précédente pour $K=1$, c'est-à-dire pour $h(t) = e^{-2\complexj\pi f_1t}$ (une seule fréquence) ? En une phrase, que fait alors $h$ au signal échantillonné $x$ ?
%         \begin{solution}
%             Avec $K=1$, $y(t) = X(-f_1)\,e^{-2\complexj\pi f_1t}$ : convoler $x$ par une unique exponentielle complexe $h(t)=e^{-2\complexj\pi f_1t}$ extrait donc, à elle seule, la valeur du spectre de $x$ à la fréquence $-f_1$. $h$ agit comme un filtre qui isole une seule composante fréquentielle du signal échantillonné.
%         \end{solution}
%     \end{enumerate}
% \end{exercice}


%% ── Ex 3 : APPLY ─────────────────────────────────────────────────
\begin{exercice}[Modulation d'un signal à bande limitée]
    Soit $x$ un signal (une fonction du temps $t\mapsto x(t)$) à bande limitée, c'est-à-dire dont la transformée de Fourier $X$ est nulle en dehors de $[-B,B]$ ($B>0$).
    
    Pour $f_0>B$, on pose
    $$
        m(t) = x(t)\sin(2\pi f_0t).
    $$
    \begin{enumerate}[\bfseries 1.]
        \item Exprimer $M=\mathcal F[m]$ en fonction de $X$.
        \begin{solution}
            Avec la formule d'Euler $\sin(2\pi f_0t) = \dfrac{e^{2\complexj\pi f_0t}-e^{-2\complexj\pi f_0t}}{2\complexj}$ et le décalage fréquentiel (vu en cours) :
            $$
                M(f) = \frac1{2\complexj}\Big(X(f-f_0) - X(f+f_0)\Big).
            $$
        \end{solution}

        \item En déduire le support de $M$, c'est-à-dire les valeurs de $f\in\R$ pour lesquelles $M(f)\neq0$.
        \begin{solution}
            $X$ est nulle en dehors de $[-B,B]$, donc $X(f-f_0)$ est nulle en dehors de $[f_0-B,f_0+B]$, et $X(f+f_0)$ en dehors de $[-f_0-B,-f_0+B]$. Comme $f_0>B$, ces deux intervalles sont disjoints, et le support de $M$ est leur réunion :
            $$
                [-f_0-B,\,-f_0+B]\ \cup\ [f_0-B,\,f_0+B].
            $$
            Moduler $x$ par un sinus translate donc son spectre autour de $\pm f_0$, sans en changer la forme.
        \end{solution}

        \item On remultiplie $m$ par le même sinus : $p(t) = m(t)\sin(2\pi f_0t)$. Exprimer $p(t)$ comme une somme de deux termes (transformation produit $\to$ somme).
        \begin{solution}
            $\sin^2(2\pi f_0t)$ se calcule avec la formule d'Euler :
            $$
                \sin^2(2\pi f_0t) = \left(\frac{e^{2\complexj\pi f_0t}-e^{-2\complexj\pi f_0t}}{2\complexj}\right)^{\!2} = \frac{e^{4\complexj\pi f_0t} - 2 + e^{-4\complexj\pi f_0t}}{-4} = \frac{1-\cos(4\pi f_0t)}2.
            $$
            D'où
            $$
                p(t) = x(t)\sin^2(2\pi f_0t) = \frac{x(t)}2 - \frac{x(t)\cos(4\pi f_0t)}2.
            $$
        \end{solution}

        \item Calculer $P=\mathcal F[p]$ en fonction de $X$.
        \begin{solution}
            Le premier terme a pour TF $\frac12X(f)$ par linéarité. Pour le second, avec la formule d'Euler $\cos(4\pi f_0t) = \frac12e^{4\complexj\pi f_0t}+\frac12e^{-4\complexj\pi f_0t}$ et le décalage fréquentiel (vu en cours, avec la fréquence $2f_0$ au lieu de $f_0$) :
            $$
                \mathcal F\big[x(t)\cos(4\pi f_0t)\big](f) = \frac12\Big(X(f-2f_0)+X(f+2f_0)\Big).
            $$
            Par linéarité,
            $$
                P(f) = \frac12X(f) - \frac14\Big(X(f-2f_0)+X(f+2f_0)\Big).
            $$
        \end{solution}

        \item Comment peut-on retrouver $x(t)$ à partir de $p(t)$ ? (Aucun calcul supplémentaire n'est demandé : justifier le principe à l'aide des questions précédentes.)
        \begin{solution}
            Comme pour $M$ (question 2), $f_0>B$ sépare les trois morceaux du spectre de $p$ : $\frac12X(f)$ occupe $[-B,B]$, tandis que les deux autres termes occupent $[2f_0-B,2f_0+B]$ et $[-2f_0-B,-2f_0+B]$, disjoints de $[-B,B]$ dès que $f_0>B$. Sur $[-B,B]$, on a donc exactement $P(f) = \frac12X(f)$.

            Un \highlight{filtre passe-bas idéal} qui laisse inchangées les fréquences de $[-B,B]$ et annule les autres, appliqué à $p$, donne donc $\frac12X(f)$ sur tout $\R$, soit $\frac12x(t)$ en temporel : on retrouve $x$ (à un facteur $2$ près) en multipliant par $2$ le résultat de ce filtrage. Moduler puis démoduler (remultiplier par la porteuse, puis filtrer passe-bas) permet ainsi de transmettre $x$ autour de $f_0$ et de le récupérer à l'arrivée.
        \end{solution}
    \end{enumerate}
\end{exercice}


%% ── Ex 4 : APPLY ─────────────────────────────────────────────────
\begin{exercice}[Convolution : deux portes]
    On note $\Pi_\theta$ la fonction porte de largeur $\theta>0$, centrée en $0$ :
    $$
        \Pi_\theta(t) = \begin{cases} 1 & \text{si } |t|<\theta/2 \\ 0 & \text{sinon.} \end{cases}
    $$
    \begin{enumerate}[\bfseries 1.]
        \item Tracer $\Pi_\theta$. Que vaut $\displaystyle\int_{-\infty}^{+\infty}\Pi_\theta(t)\,\d t$ ?
        \begin{solution}
            Un rectangle de largeur $\theta$ et de hauteur $1$. Son intégrale est l'aire de ce rectangle : $\theta$.
        \end{solution}

        \item Calculer $\mathcal F[\Pi_\theta]$.
        \begin{solution}
            $$
                \mathcal F[\Pi_\theta](f) = \int_{-\theta/2}^{\theta/2} e^{-2\complexj\pi ft}\,\d t
                = \left[\frac{e^{-2\complexj\pi ft}}{-2\complexj\pi f}\right]_{-\theta/2}^{\theta/2}
                = \frac{e^{\complexj\pi f\theta}-e^{-\complexj\pi f\theta}}{2\complexj\pi f}
                = \frac{\sin(\pi f\theta)}{\pi f} = \theta\,\sinc(\pi f\theta).
            $$
            En $f=0$, $\sinc(0)=1$ donc $\mathcal F[\Pi_\theta](0)=\theta$ : cohérent avec la question 1.
        \end{solution}

        \item Soit $x$ une fonction quelconque. Calculer $(x*\Pi_\theta)(t)$ et l'interpréter.
        \begin{solution}
            $\Pi_\theta(t-\tau)=1$ si et seulement si $|t-\tau|<\theta/2$, soit $\tau\in(t-\theta/2,\,t+\theta/2)$, et vaut $0$ sinon. Donc
            $$
                (x*\Pi_\theta)(t) = \int_{-\infty}^{+\infty} x(\tau)\,\Pi_\theta(t-\tau)\,\d\tau = \int_{t-\theta/2}^{t+\theta/2} x(\tau)\,\d\tau.
            $$
            C'est l'intégrale de $x$ sur une fenêtre de largeur $\theta$ centrée en $t$ : une aire glissante (une moyenne glissante, à un facteur $1/\theta$ près).
        \end{solution}

        \item Cas particulier $x=\Pi_\theta$ : calculer $(\Pi_\theta*\Pi_\theta)(t)$ et tracer le résultat.
        \begin{solution}
            D'après la question précédente,
            $$
                (\Pi_\theta*\Pi_\theta)(t) = \int_{t-\theta/2}^{t+\theta/2} \Pi_\theta(\tau)\,\d\tau,
            $$
            c'est-à-dire la longueur du recouvrement entre les intervalles $[-\theta/2,\theta/2]$ et $[t-\theta/2,t+\theta/2]$. Pour $|t|\geq\theta$, ces intervalles ne se recouvrent pas : le résultat est $0$. Pour $|t|<\theta$, le recouvrement vaut $\theta-|t|$. D'où
            $$
                (\Pi_\theta*\Pi_\theta)(t) = \begin{cases} \theta-|t| & \text{si } |t|<\theta \\ 0 & \text{sinon,} \end{cases}
            $$
            un triangle de base $2\theta$ et de hauteur $\theta$.
        \end{solution}

        \item En notant $Y=\mathcal F[\Pi_\theta*\Pi_\theta]$, le théorème produit-convolution donne $Y(f) = \mathcal F[\Pi_\theta](f)^2$. Sans calculer $Y$ entièrement, vérifier que $Y(0)$ est cohérent avec l'aire du triangle tracé à la question 4.
        \begin{solution}
            D'après la question 2, $\mathcal F[\Pi_\theta](0)=\theta$, donc $Y(0)=\theta^2$. Or $Y(0) = \displaystyle\int_{-\infty}^{+\infty}(\Pi_\theta*\Pi_\theta)(t)\,\d t$ est aussi l'aire du triangle de la question 4, qui vaut $\frac12\times(2\theta)\times\theta=\theta^2$. Les deux valeurs coïncident.
        \end{solution}
    \end{enumerate}
\end{exercice}


%% ── Ex 5 : APPLY ─────────────────────────────────────────────────
\begin{exercice}[Décomposition en portes décalées]
    On reprend la porte $\Pi_\theta$ définie dans l'exercice précédent, de transformée
    $\mathcal F[\Pi_\theta](f)=\theta\,\sinc(\pi f\theta)$.
    \begin{enumerate}[\bfseries 1.]
        \item Soit $x$ la fonction définie par $x(t)=A_1$ sur $[0,\theta[$, $A_2$ sur $[\theta,2\theta[$, $A_3$ sur $[2\theta,3\theta[$, et $0$ ailleurs, avec $A_1=1$, $A_2=2$, $A_3=-1$. Tracer son graphe et l'exprimer comme une somme de trois portes $\Pi_\theta$ centrées et décalées.
        \begin{solution}
            Chaque palier, de largeur $\theta$, est une porte $\Pi_\theta$ décalée pour être centrée en son milieu : le premier en $\theta/2$, le deuxième en $3\theta/2$, le troisième en $5\theta/2$. D'où
            $$
                f(t) = \Pi_\theta(t-\theta/2) + 2\,\Pi_\theta(t-3\theta/2) -\Pi_\theta(t-5\theta/2).
            $$
        \end{solution}

        \item En déduire $\mathcal F\{x\} (f)$ en reprenant les propriétés de la transformée de Fourier (linéarité, décalage).
        \begin{solution}
            Par linéarité, et en décalant $\Pi_\theta$ de $t_0$ (ce qui multiplie sa transformée par $e^{-2\complexj\pi ft_0}$) :
            $$
                \mathcal F[x](f) = \theta\,\sinc(\pi f\theta)\left(e^{-\complexj\pi f\theta} + 2\,e^{-3\complexj\pi f\theta} + e^{-5\complexj\pi f\theta}\right).
            $$
        \end{solution}

        \item Montrer que $|\mathcal F\{x\}(f)|=0$ en $f=k/\theta$ ($k\in\mathbb Z^*$), quels que soient $A_1,A_2,A_3$ et quelles que soient les positions des trois paliers.
        \begin{solution}
            D'après la question 2, $\mathcal F[x](f)$ s'écrit comme le produit de $\theta\,\sinc(\pi f\theta)$ par une somme de termes de la forme $A_k\,e^{-2\complexj\pi ft_k}$. Le premier facteur ne dépend ni des $A_k$ ni des positions $t_k$ : c'est le même dès que les trois portes ont la largeur $\theta$. Or un produit est nul dès que l'un de ses facteurs l'est : $\theta\,\sinc(\pi f\theta)=0$ pour $f=k/\theta$ ($k\in\mathbb Z^*$), donc $\mathcal F[x](f)=0$ en ces points, quelle que soit la valeur du second facteur.
        \end{solution}

        \item On reprend $(A_1,A_2,A_3)=(1,2,1)$, mais avec des largeurs différentes : $\theta_1=\theta$, $\theta_2=2\theta$, $\theta_3=3\theta$ (paliers toujours consécutifs). Exprimer $x$, puis $\mathcal F\{x\}$. L'argument de la question 3 s'applique-t-il encore ? Justifier.
        \begin{solution}
            Les paliers occupent maintenant $[0,\theta)$, $[\theta,3\theta)$, $[3\theta,6\theta)$, centrés respectivement en $\theta/2$, $2\theta$, $4{,}5\,\theta$ :
            $$
                x(t) = \Pi_\theta(t-\theta/2) + 2\,\Pi_{2\theta}(t-2\theta) + \Pi_{3\theta}(t-4{,}5\,\theta),
            $$
            $$
                \mathcal F[x](f) = \theta\,\sinc(\pi f\theta)\,e^{-\complexj\pi f\theta} + 4\theta\,\sinc(2\pi f\theta)\,e^{-4\complexj\pi f\theta} + 3\theta\,\sinc(3\pi f\theta)\,e^{-9\complexj\pi f\theta}.
            $$
            Non : chaque terme a maintenant sa propre largeur, donc son propre sinus cardinal. Il n'y a plus de facteur commun à mettre en évidence comme à la question 3 : l'argument de factorisation ne s'applique plus. Un zéro de $\mathcal F[x]$ résulterait ici d'une annulation particulière entre trois termes différents, pas d'une propriété structurelle garantie.
        \end{solution}
    \end{enumerate}
\end{exercice}


%% ── Ex 6 (optionnel) : EXTEND ─────────────────────────────────────
% \begin{exercice}[Peigne de Dirac --- approfondissement]
%     Pour $T>0$, on définit le peigne de Dirac $\Sha_T(t) = \displaystyle\sum_{n=-\infty}^{+\infty}\delta(t-nT)$.
%     \begin{enumerate}[\bfseries 1.]
%         \item Justifier que $\Sha_T$ est $T$-périodique.
%         \begin{solution}
%             $$
%                 \Sha_T(t+T) = \sum_{n=-\infty}^{+\infty}\delta(t+T-nT) = \sum_{n=-\infty}^{+\infty}\delta\big(t-(n-1)T\big).
%             $$
%             En posant $m=n-1$, qui parcourt aussi $\mathbb Z$ tout entier quand $n$ le parcourt, cette somme vaut $\displaystyle\sum_{m=-\infty}^{+\infty}\delta(t-mT) = \Sha_T(t)$.
%         \end{solution}

%         \item $\Sha_T$ admet donc une décomposition en série de Fourier complexe, de coefficients
%         $c_n=\dfrac1T\displaystyle\int_{-T/2}^{T/2}\Sha_T(t)\,e^{-\complexj n\omega_0t}\,\d t$ ($\omega_0=2\pi/T$). Sur $[-T/2,T/2)$, un seul terme de la somme définissant $\Sha_T$ contribue à l'intégrale : lequel ? En déduire $c_n$. On rappelle que $\displaystyle\int_{-\infty}^{+\infty} x(t)\,\delta(t)\,\d t = x(0)$ pour toute fonction $x$.
%         \begin{solution}
%             Seul le terme $n=0$, $\delta(t)$, a son support (le point $t=0$) dans $[-T/2,T/2)$ ; tout autre terme $\delta(t-kT)$ ($k\neq0$) est concentré en $t=kT$, hors de cet intervalle. Donc
%             $$
%                 c_n = \frac1T\int_{-T/2}^{T/2}\delta(t)\,e^{-\complexj n\omega_0t}\,\d t = \frac1T\,e^{-\complexj n\omega_0\times0} = \frac1T, \qquad \text{pour tout } n\in\mathbb Z.
%             $$
%         \end{solution}

%         \item Écrire la série de Fourier complexe de $\Sha_T(t)$ obtenue à la question précédente.
%         \begin{solution}
%             $$
%                 \Sha_T(t) = \sum_{n=-\infty}^{+\infty} c_n\,e^{\complexj n\omega_0t} = \frac1T\sum_{n=-\infty}^{+\infty} e^{\complexj n\omega_0t}.
%             $$
%         \end{solution}

%         \item On rappelle que $\mathcal F\big[e^{2\complexj\pi f_0t}\big](f) = \delta(f-f_0)$ pour tout $f_0\in\R$. En l'appliquant terme à terme à la question précédente, calculer $\mathcal F[\Sha_T]$.
%         \begin{solution}
%             Chaque terme $e^{\complexj n\omega_0t} = e^{2\complexj\pi(n/T)t}$ a pour transformée $\delta(f-n/T)$. Par linéarité,
%             $$
%                 \mathcal F[\Sha_T](f) = \frac1T\sum_{n=-\infty}^{+\infty} \delta(f-n/T).
%             $$
%             C'est de nouveau un peigne de Dirac, d'espacement $1/T$ et de hauteur $1/T$.
%         \end{solution}

%         \item Soit $x(t)$ une fonction quelconque et $y(t)=x(t)\cdot\Sha_T(t)$. Le théorème produit-convolution énonce que $\mathcal F[uv]=\mathcal F[u]*\mathcal F[v]$ pour tous signaux $u,v$. En déduire $\mathcal F[y]$ en fonction de $X=\mathcal F[x]$.
%         \begin{solution}
%             Avec $u=x$, $v=\Sha_T$ :
%             $$
%                 \mathcal F[y](f) = X(f) * \frac1T\sum_{n=-\infty}^{+\infty}\delta(f-n/T).
%             $$
%         \end{solution}

%         \item On rappelle que $X*\delta(\cdot-f_0) = X(\cdot-f_0)$ pour toute fonction $X$ et tout $f_0\in\R$. En déduire une expression simplifiée de $Y=\mathcal F[y]$.
%         \begin{solution}
%             Par linéarité de la convolution,
%             $$
%                 Y(f) = \frac1T\sum_{n=-\infty}^{+\infty} X(f)*\delta(\cdot-n/T) = \frac1T\sum_{n=-\infty}^{+\infty} X(f-n/T).
%             $$
%         \end{solution}

%         \item Montrer, en réindexant la somme, que $Y(f+1/T)=Y(f)$ pour tout $f$ : $Y$ est $1/T$-périodique. Cette périodicité dépend-elle d'une périodicité de $X$ ?
%         \begin{solution}
%             $$
%                 Y(f+1/T) = \frac1T\sum_{n=-\infty}^{+\infty} X\big(f+1/T-n/T\big) = \frac1T\sum_{n=-\infty}^{+\infty} X\big(f-(n-1)/T\big).
%             $$
%             En posant $m=n-1$, qui parcourt aussi $\mathbb Z$ tout entier quand $n$ le parcourt :
%             $$
%                 Y(f+1/T) = \frac1T\sum_{m=-\infty}^{+\infty} X(f-m/T) = Y(f).
%             $$
%             Cette égalité ne repose sur aucune propriété de $x$ : $x$ était une fonction quelconque dès le départ. La périodicité de $Y$ vient donc entièrement de la structure du peigne $\Sha_T$, et non d'une éventuelle périodicité de $X$.
%         \end{solution}
%     \end{enumerate}
% \end{exercice}

\end{document}
